\begin{figure}[H]
\centering

\begin{center}
\textbf{Completeness Dependency Map}
\end{center}

\begin{tikzpicture}[>=Latex]

% Nodes
\node[depaxiom]   (comp)  at (0,0)       {Completeness Axiom};

\node[deptheorem] (nip)   at (-2.9,-1.7) {Nested Interval\\Property};
\node[deptheorem] (sqrt)  at ( 3.4,-1.7) {Existence of $\sqrt{a}$};

\node[deptheorem] (arch)  at (-2.9,-3.3) {Archimedean Property};
\node[deptheorem] (floor) at (-2.9,-4.9) {Floor Lemma};
\node[deptheorem] (densQ) at (-2.9,-6.5) {Density of $\mathbb{Q}$};
\node[deptheorem] (densI) at (-2.9,-8.1) {Density of the\\Irrationals in $\mathbb{R}$};

% Edges
\draw[depimplies] (comp) -- (nip);
\draw[depimplies] (comp) -- (sqrt);
\draw[depimplies] (nip)  -- (arch);
\draw[depimplies] (arch) -- (floor);
\draw[depimplies] (floor) -- (densQ);
\draw[depimplies] (densQ) -- (densI);

\end{tikzpicture}

\begin{center}
\DependencyFigureLegend
\end{center}

\caption{Dependency map for the completeness subchapter. The Completeness Axiom yields the Nested Interval Property and the existence of $\sqrt{a}$; the Nested Interval Property then supports the Archimedean Property, the Floor Lemma, and the density results for $\mathbb{Q}$ and for the irrational numbers.}
\label{fig:completeness-chain}
\end{figure}